\section{Zadanie 4: Obliczanie n-tej sumy i badanie zbieżności szeregów}
\label{sec:zadanie4}

\begin{tresczadania}
Obliczyć $n$-tą sumę i zbadać zbieżność szeregów:
\begin{enumerate}[label=(\alph*),leftmargin=2em]
    \item $\displaystyle \sum_{n=1}^{\infty} \ln\left(1 + \frac{1}{n}\right)$
    \item $\displaystyle 1 + \sum_{n=2}^{\infty} \ln\left(1 - \frac{1}{n^2}\right)$
\end{enumerate}
\end{tresczadania}

\vspace{0.2cm}

\begin{teoriabox}
\begin{enumerate}
    \item \textbf{Ciąg sum częściowych ($n$-ta suma):}
    Dla szeregu liczb rzeczywistych $\sum_{k=1}^{\infty} a_k$, ciąg sum częściowych definiuje się jako:
    \begin{equation}
        S_n \coloneqq \sum_{k=1}^{n} a_k = a_1 + a_2 + \dots + a_n.
        \label{eq:partial_sum_def}
    \end{equation}

    \item \textbf{Definicja zbieżności szeregu:}
    Szereg jest \textbf{zbieżny} wtedy i tylko wtedy, gdy istnieje skończona granica ciągu sum częściowych:
    \begin{equation*}
        S = \lim_{n \to \infty} S_n \in \mathbb{R}.
    \end{equation*}
    Wartość $S$ nazywamy \textbf{sumą szeregu}. Gdy granica jest nieskończona lub nie istnieje, szereg jest \textbf{rozbieżny}.

    \item \textbf{Własności logarytmu naturalnego i iloczynów:}
    Dla liczb dodatnich zachodzą prawa:
    \begin{equation*}
        \ln(x \cdot y) = \ln x + \ln y, \qquad \ln\left(\frac{x}{y}\right) = \ln x - \ln y \implies \sum_{k=1}^{n} \ln(u_k) = \ln\left( \prod_{k=1}^{n} u_k \right).
    \end{equation*}

    \item \textbf{Ciągłość funkcji logarytmicznej:}
    Z uwagi na ciągłość funkcji $\ln$ na $(0, \infty)$, przejście graniczne można przenieść do argumentu logarytmu: $\lim \ln(u_n) = \ln(\lim u_n)$, o ile $\lim u_n > 0$.
\end{enumerate}
\end{teoriabox}

\vspace{0.2cm}

\subsection*{Rozwiązania szczegółowe}

% -----------------------------------------------------------------------------
\subsubsection*{Podpunkt (a): $\displaystyle \sum_{n=1}^{\infty} \ln\left(1 + \frac{1}{n}\right)$}

\begin{enumerate}[label=\textbf{Krok \arabic*.},leftmargin=*]
    \item \textbf{Przekształcenie wyrazu ogólnego.}
    Zapisujemy argument pod wspólnym mianownikiem i rozbijamy logarytm ilorazu:
    \begin{equation*}
        a_k = \ln\left(1 + \frac{1}{k}\right) = \ln\left(\frac{k+1}{k}\right) = \ln(k+1) - \ln k.
    \end{equation*}

    \item \textbf{Wyznaczenie $n$-tej sumy częściowej $S_n$.}
    Zapisujemy sumę wyrazów od $k = 1$ do $n$:
    \begin{align*}
        S_n &= \sum_{k=1}^{n} \big( \ln(k+1) - \ln k \big) \\
        &= (\ln 2 - \ln 1) + (\ln 3 - \ln 2) + (\ln 4 - \ln 3) + \dots + \big( \ln(n+1) - \ln n \big).
    \end{align*}
    Wyrazy pośrednie redukują się wzajemnie (zwijanie teleskopowe):
    \begin{equation*}
        S_n = \ln(n+1) - \ln 1 = \ln(n+1) - 0 = \ln(n+1).
    \end{equation*}
    (Iloczynowo: $S_n = \ln\left(\prod_{k=1}^n \frac{k+1}{k}\right) = \ln\left(\frac{2}{1} \cdot \frac{3}{2} \cdots \frac{n+1}{n}\right) = \ln(n+1)$).

    \item \textbf{Badanie zbieżności szeregu.}
    Obliczamy granicę sum częściowych przy $n \to \infty$:
    \begin{equation*}
        \lim_{n \to \infty} S_n = \lim_{n \to \infty} \ln(n+1) = +\infty.
    \end{equation*}

    \item \textbf{Wniosek.}
    Ponieważ granica ciągu sum częściowych wynosi $+\infty$, szereg $\displaystyle\sum_{n=1}^{\infty} \ln\left(1 + \frac{1}{n}\right)$ jest \textbf{rozbieżny}.
\end{enumerate}

% -----------------------------------------------------------------------------
\subsubsection*{Podpunkt (b): $\displaystyle 1 + \sum_{n=2}^{\infty} \ln\left(1 - \frac{1}{n^2}\right)$}

\begin{enumerate}[label=\textbf{Krok \arabic*.},leftmargin=*]
    \item \textbf{Faktoryzacja argumentu logarytmu.}
    Dla $k \geqslant 2$ stosujemy wzór na różnicę kwadratów:
    \begin{equation*}
        1 - \frac{1}{k^2} = \frac{k^2 - 1}{k^2} = \frac{(k-1)(k+1)}{k^2} = \left(\frac{k-1}{k}\right) \cdot \left(\frac{k+1}{k}\right).
    \end{equation*}

    \item \textbf{Obliczenie sumy składników szeregu (od $k = 2$ do $n$).}
    Oznaczmy sumę częściową samej sumy szeregu przez $T_n$:
    \begin{equation*}
        T_n = \sum_{k=2}^{n} \ln\left(1 - \frac{1}{k^2}\right) = \ln\left( \prod_{k=2}^{n} \frac{(k-1)(k+1)}{k^2} \right).
    \end{equation*}
    Rozbijamy iloczyn na dwa niezależne iloczyny teleskopowe $P_n$ i $Q_n$:
    \begin{align*}
        P_n &= \prod_{k=2}^{n} \frac{k-1}{k} = \frac{1}{2} \cdot \frac{2}{3} \cdot \frac{3}{4} \cdots \frac{n-1}{n} = \frac{1}{n}, \\
        Q_n &= \prod_{k=2}^{n} \frac{k+1}{k} = \frac{3}{2} \cdot \frac{4}{3} \cdot \frac{5}{4} \cdots \frac{n+1}{n} = \frac{n+1}{2}.
    \end{align*}
    Iloczyn całkowity wynosi zatem:
    \begin{equation*}
        \prod_{k=2}^{n} \frac{(k-1)(k+1)}{k^2} = P_n \cdot Q_n = \frac{1}{n} \cdot \frac{n+1}{2} = \frac{n+1}{2n}.
    \end{equation*}
    Stąd:
    \begin{equation*}
        T_n = \ln\left(\frac{n+1}{2n}\right).
    \end{equation*}

    \item \textbf{Wyznaczenie $n$-tej sumy całego wyrażenia.}
    Całe zadane wyrażenie zawiera początkowy wyraz $1$, stąd pełna $n$-ta suma wynosi:
    \begin{equation*}
        S_n = 1 + T_n = 1 + \ln\left(\frac{n+1}{2n}\right).
    \end{equation*}

    \item \textbf{Badanie zbieżności i obliczenie sumy szeregu.}
    Obliczamy granicę przy $n \to \infty$:
    \begin{equation*}
        \lim_{n \to \infty} \frac{n+1}{2n} = \lim_{n \to \infty} \left(\frac{1}{2} + \frac{1}{2n}\right) = \frac{1}{2}.
    \end{equation*}
    Z ciągłości logarytmu:
    \begin{equation*}
        \lim_{n \to \infty} \ln\left(\frac{n+1}{2n}\right) = \ln\left(\frac{1}{2}\right) = -\ln 2.
    \end{equation*}
    Zatem granica sumy wynosi:
    \begin{equation*}
        S = \lim_{n \to \infty} S_n = 1 + \ln\left(\frac{1}{2}\right) = 1 - \ln 2.
    \end{equation*}

    \item \textbf{Wniosek.}
    Granica $S$ istnieje i jest skończona ($S = 1 - \ln 2 \approx 0{,}3069$). Zatem szereg jest \textbf{zbieżny}, a jego suma wynosi $1 - \ln 2$.
\end{enumerate}

\vspace{0.2cm}

\begin{odpowiedzbox}[Podsumowanie wyników Zadania 4]
\begin{enumerate}[label=(\alph*),leftmargin=2em]
    \item Szereg $\displaystyle\sum_{n=1}^{\infty} \ln\left(1 + \frac{1}{n}\right)$:
    \begin{itemize}
        \item $n$-ta suma częściowa: $\mathbf{S_n = \ln(n+1)}$,
        \item Zbieżność: szereg jest \textbf{rozbieżny} ($\lim_{n\to\infty} S_n = +\infty$).
    \end{itemize}
    \item Wyrażenie $\displaystyle 1 + \sum_{n=2}^{\infty} \ln\left(1 - \frac{1}{n^2}\right)$:
    \begin{itemize}
        \item $n$-ta suma: $\mathbf{S_n = 1 + \ln\left(\frac{n+1}{2n}\right)}$,
        \item Zbieżność: szereg jest \textbf{zbieżny}, a jego suma wynosi $\mathbf{S = 1 - \ln 2}$.
    \end{itemize}
\end{enumerate}
\end{odpowiedzbox}
